\subsection{ALGEBRAIC OPERATIONS ON REAL-VALUED FUNCTIONS}

Let \(f\) and \(g\) be two real-valued functions defined on a set \(\mathbf{X}\) ; if the sum \(f(x) + g(x)\) {[}resp. the product \(f(x)g(x)\) {]} is defined for all \(x \in \mathbf{X}\) , then we denote by \(f + g\) (resp. \(fg\) ) the real-valued function

\[
x \rightarrow f (x) + g (x) \quad [ \text { resp. } x \rightarrow f (x) g (x) ].
\]

Again, if \(1 / f(x)\) is defined for all \(x\in \mathbf{X}\) , then \(1 / f\) denotes the function \(x\to 1 / f(x)\) .

This last function is therefore defined provided \(f\) does not take the value \(0\) ; when \(f\) takes its values in the interval \([0, +\infty]\) (resp. in \([-\infty, 0]\) ) \(1/f(x)\) as everywhere defined by putting \(1/0 = +\infty\) (resp. \(1/0 = -\infty\) ); in this case the function \(1/f\) is defined.

Suppose that X is filtered by a filter \(\mathfrak{F}\) , and that \(\lim_{\mathfrak{F}} f\) and \(\lim_{\mathfrak{F}} g\) exist. If on the one hand the function \(f + g\) (resp. \(fg\) , \(1/g\) ) is defined, and if on the other hand the expression \(\lim_{\mathfrak{F}} f + \lim_{\mathfrak{F}} g\) (resp. \(\lim_{\mathfrak{F}} f\) . \(\lim_{\mathfrak{F}} g\) , \(1/\lim_{\mathfrak{F}} f\) ) has a sense, then \(\lim_{\mathfrak{F}} (f + g)\) {[}resp. \(\lim_{\mathfrak{F}} fg\) , \(\lim_{\mathfrak{F}} (1/f)\) {]} exists and is equal to this expression by reason of the continuity of the function \(x + y\) (resp. \(xy\) , \(1/x\) ) at points where it is defined.

PROPOSITION 12. Let \(f\) and \(g\) be two real-valued functions defined on a set \(\mathbf{X}\) , and let \(\mathbf{A}\) be a non-empty subset of \(\mathbf{X}\) .

\textit{(i) We have}

\begin{enumerate}
\def\labelenumi{(\arabic{enumi})}
\setcounter{enumi}{16}
\tightlist
\item
\end{enumerate}

\[
\sup _ {x \in \mathbf {A}} (f (x) + g (x)) \leqslant \sup _ {x \in \mathbf {A}} f (x) + \sup _ {x \in \mathbf {A}} g (x),\tag{17}
\]

\[
\sup _ {x \in \mathbf {A}} f (x) + \inf _ {x \in \mathbf {A}} g (x) \leqslant \sup _ {x \in \mathbf {A}} (f (x) + g (x)),\tag{18}
\]

whenever both sides of these inequalities are defined.

\begin{enumerate}
\def\labelenumi{(\roman{enumi})}
\setcounter{enumi}{1}
\tightlist
\item
  If \(f(x)\) and \(g(x)\) are \(\geqslant 0\) for each \(x \in \mathbf{A}\) , then
\end{enumerate}

\begin{enumerate}
\def\labelenumi{(\arabic{enumi})}
\setcounter{enumi}{18}
\tightlist
\item
\end{enumerate}

\[
\sup _ {x \in \mathbf {A}} (f (x) g (x)) \leqslant \sup _ {x \in \mathbf {A}} f (x) \sup _ {x \in \mathbf {A}} g (x),\tag{19}
\]

\[
\sup _ {x \in \mathbf {A}} f (x) \inf _ {x \in \mathbf {A}} g (x) \leqslant \sup _ {x \in \mathbf {A}} (f (x) g (x))\tag{20}
\]

whenever both sides of these inequalities are defined.

\begin{enumerate}
\def\labelenumi{(\roman{enumi})}
\setcounter{enumi}{2}
\tightlist
\item
  If \(f(x) \geqslant 0\) for all \(x \in \mathbf{A}\) , then
\end{enumerate}

\[
\sup _ {x \in \mathbf {A}} (1 / f (x)) = \inf _ {x \in \mathbf {A}} f (x)\tag{21}
\]

(putting \(1 / 0 = +\infty\) ).

Let \(\mathbf{H}\) be any finite subset of \(\mathbf{A}\) . If \(x_0\) is one of the points of \(\mathbf{H}\) where \(f + g\) takes its greatest value, then we have

\[
f (x _ {0}) + g (x _ {0}) \leqslant \sup _ {x \in \mathbf {H}} f (x) + \sup _ {x \in \mathbf {H}} g (x);
\]

on the other hand, if \(x_{1}\) is one of the points of \(H\) where \(f\) takes its greatest value, then

\[
f (x _ {1}) + g (x _ {1}) \geqslant \sup _ {x \in \mathbf {H}} f (x) + \inf _ {x \in \mathbf {H}} g (x);
\]

therefore

\[
\sup _ {x \in \mathbf {H}} f (x) + \inf _ {x \in \mathbf {H}} g (x) \leqslant \sup _ {x \in \mathbf {H}} (f (x) + g (x)) \leqslant \sup _ {x \in \mathbf {H}} f (x) + \sup _ {x \in \mathbf {H}} g (x).
\]

The inequalities (17) and (18) follow from this by applying Proposition 5 of no. 4 and Theorem 1 of no. 2. The proofs of the other inequalities are analogous.

COROLLARY 1. Let \(f\) be a real-valued function defined on \(\mathbf{X}\) , and let \(k\) be a real number. Then

\[
\sup _ {x \in \mathbf {A}} (f (x) + k) = k + \sup _ {x \in \mathbf {A}} f (x)\tag{22}
\]

whenever both sides are defined, and, if \(k \geqslant 0\) ,

\[
\sup _ {x \in \mathbf {A}} (k f (x)) = k. \sup _ {x \in \mathbf {A}} f (x)\tag{23}
\]

whenever both sides are defined.

COROLLARY 2. Let \(f_{1}\) and \(f_{2}\) be two real-valued functions defined on sets \(X_{1}, X_{2}\) respectively; then if \(A_{1}, A_{2}\) are any non-empty subsets of \(X_{1}, X_{2}\) respectively, we have

\[
\sup _ {(x _ {1}, x _ {2}) \in \mathbf {A} _ {1} \times \mathbf {A} _ {2}} (f _ {1} (x _ {1}) + f _ {2} (x _ {2})) = \sup _ {x _ {1} \in \mathbf {A} _ {1}} f _ {1} (x _ {1}) + \sup _ {x _ {2} \in \mathbf {A} _ {2}} f _ {2} (x _ {2})\tag{24}
\]

whenever both sides are defined; and if \(f_1, f_2\) are \(\geqslant 0\) in \(A_1, A_2\) respectively we have

\[
\sup _ {(x _ {1}, x _ {2}) \in \mathbf {A} _ {1} \times \mathbf {A} _ {2}} (f _ {1} (x _ {1}) f _ {2} (x _ {2})) = \sup _ {x _ {1} \in \mathbf {A} _ {1}} f _ {1} (x _ {1}) \sup _ {x _ {2} \in \mathbf {A} _ {2}} f _ {2} (x _ {2}),\tag{25}
\]

whenever both sides are defined.

This is an immediate consequence of the preceding corollary and Proposition 9 of no. 4.

In particular, if \(\mathbf{A}\) and \(\mathbf{B}\) are two subsets of \(\overline{\mathbf{R}}\) such that the set \(\mathbf{A} + \mathbf{B}\) of sums \(x + y\) ( \(x \in \mathbf{A}, y \in \mathbf{B}\) ) is defined, we have

\[
\sup (\mathbf {A} + \mathbf {B}) = \sup \mathbf {A} + \sup \mathbf {B}\tag{26}
\]

if the right-hand side is defined. Again, if A and B are two subsets of \([0, +\infty]\) , we have

\[
\sup \mathbf {A B} = \sup \mathbf {A}. \sup \mathbf {B}\tag{27}
\]

whenever both sides are defined.

PROPOSITION 13. Let \(f\) and \(g\) be two real-valued functions defined on a filtered set \(X\) .

\begin{enumerate}
\def\labelenumi{(\roman{enumi})}
\tightlist
\item
  We have
\end{enumerate}

\[
\lim \sup (f + g) \leqslant \lim \sup f + \lim \sup g, \tag {28}
\]

\begin{enumerate}
\def\labelenumi{(\arabic{enumi})}
\setcounter{enumi}{28}
\tightlist
\item
  \(\lim \sup f + \lim \inf g \leqslant \lim \sup (f + g)\)
\end{enumerate}

whenever both sides of these inequalities are defined.

\begin{enumerate}
\def\labelenumi{(\roman{enumi})}
\setcounter{enumi}{1}
\tightlist
\item
  If \(f\) and \(g\) are \(\geqslant 0\) on \(\mathbf{X}\) , we have
\end{enumerate}

\[
\lim \sup f g \leqslant (\lim \sup f) (\lim \sup g), \tag {30}
\]

\[
(\lim \sup f) (\lim \inf g) \leqslant \lim \sup f g \tag {31}
\]

whenever both sides of these inequalities are defined.

\begin{enumerate}
\def\labelenumi{(\roman{enumi})}
\setcounter{enumi}{2}
\tightlist
\item
  If \(f \geqslant 0\) on \(\mathbf{X}\) , then
\end{enumerate}

\[
\lim \sup (1 / f) = 1 / (\lim \inf f)\tag{32}
\]

(putting \(1 / 0 = +\infty\) ).

These relations are consequences of Proposition 12 and relations (10)

COROLLARY 1. Let \(f\) and \(g\) be two real-valued functions defined on a filtered set \(X\) . If \(\lim g\) exists, then

\[
\lim \sup (f + g) = \lim \sup f + \lim g \tag {33}
\]

whenever both sides are defined, and

\[
\lim \sup f g = (\lim \sup f) (\lim g) \tag {34}
\]

whenever both sides are defined and \(f\) and \(g\) are \(\geqslant 0\) .

COROLLARY 2. Let \(f\) and \(g\) be two real-valued functions defined on a filtered set X. If \(\lim f = +\infty\) and \(\lim \inf g > -\infty\) and \(f + g\) is defined, then \(\lim (f + g) = +\infty\) . If \(\lim f = +\infty\) and \(\lim \inf g > 0\) and \(fg\) is defined, then \(\lim fg = +\infty\) .

